\documentclass[11pt]{article}
\usepackage[T1]{fontenc}
\usepackage{lmodern,amsmath,amssymb,amsthm,mathtools,microtype}
\usepackage[letterpaper,margin=.9in]{geometry}
\usepackage{hyperref}
\hypersetup{pdfauthor={Peyton Wang and Daniel Wang},pdftitle={Rational Pairs Suffice for Illumination in Orthogonal Lattice Rooms},colorlinks=true,linkcolor=black,citecolor=black,urlcolor=blue}
\setlength{\parindent}{0pt}
\setlength{\parskip}{4pt plus 1pt minus 1pt}
\newtheorem{theorem}{Theorem}
\newtheorem{lemma}{Lemma}
\newtheorem{proposition}{Proposition}
\newtheorem{corollary}{Corollary}
\theoremstyle{remark}
\newtheorem{remark}{Remark}
\newcommand{\T}{\mathbb T}
\newcommand{\Q}{\mathbb Q}
\newcommand{\Z}{\mathbb Z}
\newcommand{\R}{\mathbb R}
\newcommand{\interior}{\operatorname{int}}
\DeclareMathOperator{\diag}{diag}
\title{Rational Pairs Suffice for Illumination\\in Orthogonal Lattice Rooms}
\author{Peyton Wang\\\small Francis Parker School \and Daniel Wang}
\date{Working draft v0.1 --- 10 November 2026\\\small Not submitted or peer reviewed}
\begin{document}
\maketitle
\begin{abstract}
Let $P$ be a strictly simple orthogonal polygon with integer vertices.
We prove that if every pair of distinct rational interior points illuminates,
then every pair of distinct real interior points illuminates.
More precisely, any dark pair with an irrational coordinate belongs locally
to a one-parameter family of dark pairs containing rational pairs arbitrarily
nearby. The proof uses the four-copy square-tiled unfolding, horizontal and
vertical multi-twists, and elementary density arguments on circles and tori.
The reduction gives no bound on rational denominators and requires no
limiting deformation of the room.
\end{abstract}

\section{Statement and conventions}
A room is the closed region bounded by a Jordan polygonal curve.
Throughout, $P$ is orthogonal, has integer vertices, and has no holes,
self-contacts, or interior mirrors. A legal billiard path has finitely many
straight legs, reflects specularly at smooth boundary points, and never
meets a polygon vertex. Distinct interior points \emph{illuminate} one another
if a legal path connects them; otherwise the pair is \emph{dark}.
A point is rational when both coordinates are rational.

\begin{theorem}[Rational-pair reduction]\label{thm:reduction}
If every pair of distinct rational points in $\interior(P)$ illuminates,
then every pair of distinct points in $\interior(P)$ illuminates.
\end{theorem}

We prove the following stronger local statement.
\begin{theorem}[A dark irrational pair slides]\label{thm:slide}
Suppose $(S,T)$ is a dark pair of distinct interior points and $S$ is not
rational. After exchanging the coordinate axes if necessary, there are
$x\in\R\setminus\Q$, $y,v,c\in\Q$, and $\rho\in\Q\setminus\{0\}$ such that
\[
 S=(x,y),\qquad T=(\rho x+c,v).
\]
For all sufficiently small real $\delta$, the distinct interior pair
\[
 S_\delta=(x+\delta,y),\qquad
 T_\delta=(\rho x+c+\rho\delta,v)
\]
is dark. In particular, there are rational dark pairs arbitrarily close
to $(S,T)$.
\end{theorem}

The result concerns a fixed room. It does not say that dark pairs are
rational, that their number is finite, or that a bounded-denominator grid
can certify illumination. No wall-count bound is imposed.

\section{The marked four-copy unfolding}
Let $K=\{1,-1\}^2$. Form $M$ from four copies of $P$ labelled by
$e=(e_1,e_2)\in K$, with developing charts $(e_1x,e_2y)$.
Glue vertical boundary edges to the sheet with $e_1$ reversed, and horizontal
boundary edges to the sheet with $e_2$ reversed. The resulting connected
translation surface is tiled by unit squares. The map
\[
 \pi:M\longrightarrow\T^2=\R^2/\Z^2,
 \qquad \pi(S,e)=(e_1x,e_2y)\pmod{\Z^2}
\]
is a finite covering away from $0$. Put $L=\pi^{-1}(0)$ and let
$\Sigma\subset L$ consist of all lifts of polygon vertices.
The lifts of convex vertices are regular metric points, but remain
\emph{marked forbidden points}. Smooth boundary lattice points and
interior lattice points are not forbidden.

Call a pair of distinct points in $M\setminus\Sigma$ \emph{surface-lit}
if a straight geodesic segment joins them while avoiding $\Sigma$.
For distinct room points $S,T$, fix $s=(S,(1,1))$ and write $t_e=(T,e)$.
Then
\begin{equation}\label{eq:unfold}
 (S,T)\text{ illuminates}\quad\Longleftrightarrow\quad
 (s,t_e)\text{ is surface-lit for some }e\in K.
\end{equation}
Indeed, reflect each successive billiard leg into the adjacent sheet;
conversely project a geodesic through the sheet identifications.
Marking all vertex lifts makes corner avoidance identical in both directions.
A compact geodesic avoiding the finite marked set has only finitely many
crossings of the finite polygon edges. For interior endpoints, a segment
cannot run along a boundary edge. This is the usual unfolding correspondence;
see also \cite[Lemma 2.1 and Remark 2.2]{Wolecki}.

\begin{lemma}[Openness]\label{lem:open}
The surface-lit pairs form an open subset of
$(M\setminus\Sigma)^2$ minus the diagonal.
\end{lemma}
\begin{proof}
A legal segment is compact and has positive distance from $\Sigma$.
Continue it slightly at both ends and develop a sufficiently thin strip
around it by parallel transport. The strip maps locally isometrically to
$M\setminus\Sigma$; global injectivity is unnecessary.
Small changes of the two endpoints in this strip are joined by a straight
segment remaining inside the strip. Its image is a legal geodesic.
The endpoint charts are local homeomorphisms, proving openness.
\end{proof}

\subsection{Two multi-twists}
Cut $M$ along all horizontal grid edges. Each component is a horizontal
row cylinder of height $1$ and integer circumference. Choose a positive
integer $A$ divisible by every such circumference. On a cylinder with
coordinates $(u,\eta)\in(\R/\ell\Z)\times[0,1]$, set
\[
 h(u,\eta)=(u+A\eta,\eta).
\]
The boundary values are the identity, so the maps glue to an affine
homeomorphism of $M$ fixing every horizontal grid edge pointwise.
Similarly there is a vertical multi-twist $v$, with a positive integer
parameter $B$, fixing every vertical grid edge pointwise. Their derivatives
and torus actions are
\[
 Dh=\begin{pmatrix}1&A\\0&1\end{pmatrix}=H,
 \qquad Dv=\begin{pmatrix}1&0\\B&1\end{pmatrix}=V,
 \qquad \pi h=H\pi,\quad\pi v=V\pi.
\]
Both fix $L$, and hence $\Sigma$, pointwise.
Every element of $G=\langle h,v\rangle$ preserves surface illumination:
it maps straight segments to straight segments and preserves the forbidden
set. The diagonal action of $G$ on $M\times M$ will be used below.
We do not assume that this action descends to pairs of room points.

\section{Density under horizontal and vertical twists}
We use the elementary torus-rotation fact that multiples of a vector
$(\alpha_1,\ldots,\alpha_d)$ are dense modulo a period lattice when
$1,\alpha_1,\ldots,\alpha_d$ are rationally independent after normalizing
the periods. For a circle this is the irrational-rotation theorem;
for a product of circles it is Kronecker's theorem.

\begin{lemma}[One-point density]\label{lem:one}
Let $X$ be a compact connected square-tiled translation surface, with tile
side $1/N$ and tiling map to $\T^2$. Suppose affine maps $h_0,v_0$ fix all
tile vertices and have derivatives
$\left(\begin{smallmatrix}1&a\\0&1\end{smallmatrix}\right)$ and
$\left(\begin{smallmatrix}1&0\\b&1\end{smallmatrix}\right)$,
where $a,b$ are nonzero integers. Every point with an irrational torus
coordinate has dense $\langle h_0,v_0\rangle$-orbit.
\end{lemma}
\begin{proof}
First pass to positive powers that fix every horizontal, respectively
vertical, tile edge pointwise. Such powers exist: there are finitely many
edges, each is mapped to an edge of the same direction and length, and
an edge preserved with its orientation is fixed pointwise. Fixing vertices
alone would not suffice, since different prongs can be permuted.
The derivatives of these powers still have nonzero integral shear parameters.
It suffices to prove density for the resulting subgroup.

Suppose the second torus coordinate is irrational. The point lies in an
open row cylinder at an irrational transverse height. The horizontal twist
acts on its horizontal leaf by an irrational rotation, because the
circumference and tile offsets are rational. Thus the orbit closure $Y$
contains that entire leaf. It crosses each tile it meets horizontally.
In any column cylinder met by the leaf, points on the leaf with irrational
first coordinate have dense vertical-twist orbits on their vertical leaves.
These leaves are dense in that column cylinder, so its closure lies in $Y$.

Call a row or column cylinder full if its closure lies in $Y$. If a full
cylinder and one of the other direction share a tile, the second is full:
use the irrational transverse coordinates in their common tile and the
corresponding twist. Every pair of edge-adjacent tiles belongs to a common
row or column cylinder. Connectedness propagates fullness through every
tile. If only the first coordinate is irrational, interchange the axes.
\end{proof}

\begin{lemma}[Mixed rationality]\label{lem:mixed}
If $\pi(s)$ is rational and $\pi(t)$ is not, then $(s,t)$ is surface-lit.
\end{lemma}
\begin{proof}
The $G$-orbit of $s$ is finite: its torus coordinates have a fixed bounded
denominator, and $\pi$ has finite fibres. Choose $n>0$ such that both
$h^n$ and $v^n$ fix $s$. Lemma~\ref{lem:one} makes the orbit of $t$ under
$\langle h^n,v^n\rangle$ dense. Move it into a sufficiently small flat disc
around the regular point $s$, disjoint from $\Sigma$. The distinct endpoints
then join by a segment. Apply the inverse affine map.
\end{proof}

A pair $(z,w)\in\T^2\times\T^2$ is \emph{generic} here if no nonzero
$(a,b)\in\Z^2$ satisfies $az+bw\in\Q^2/\Z^2$.

\begin{lemma}[Generic-pair density]\label{lem:generic}
If $(\pi(s),\pi(t))$ is generic, the diagonal $G$-orbit of $(s,t)$ is dense
in $M\times M$.
\end{lemma}
\begin{proof}
Write $\pi(s)=(x,y)$ and $\pi(t)=(u,w)$, and let $Y$ be the orbit closure.
Let $r_x$ be the rank of $x,u$ in the vector space $\R/\Q$ over $\Q$,
and $r_y$ the rank of $y,w$.

If $r_y=2$, then $1,y,w$ are rationally independent. Horizontal twisting
fills the product of the horizontal leaves through $s,t$: their
circumferences are integers, so period normalization preserves the
independence. On this product, points with $1,x',u'$ rationally independent
are dense. Vertical twisting fills their products of vertical leaves.
Hence $Y$ contains the product of the closures of every pair of column
cylinders met by those horizontal leaves.

A pair of cylinders of the same direction is full if the product of its
closures lies in $Y$. If $(C_1,C_2)$ is full and each $C_i'$ intersects
$C_i$ in a tile, with the opposite direction, then $(C_1',C_2')$ is full:
apply the same independent-rotation argument in the product of the two
intersection tiles. This propagates to every tile pair. To move one tile
factor while holding the other, alternate horizontal and vertical cylinder
steps; in the stationary factor use its row and column through its fixed
tile. A step can be padded by staying in that tile, so both factors can
follow the same alternating schedule. Connectedness of the tile adjacency
graph therefore gives $Y=M\times M$.
The case $r_x=2$ follows by exchanging axes.

The only remaining generic case is $r_x=r_y=1$. Write
\[
 (x,u)=q'+\xi\beta,\qquad (y,w)=q+\eta\alpha,
\]
with $q,q',\alpha,\beta\in\Q^2$, $\alpha,\beta\ne0$, and $\xi,\eta$
irrational. Genericity forces $\det(\alpha,\beta)\ne0$, since otherwise
one nonzero rational row vector annihilates both irrational parts.
After taking a multiple to kill the rational offsets, horizontal twisting
has a circle orbit closure in direction $\alpha$. Consequently $Y$
contains pairs obtained by moving the first and second points horizontally
by $r\alpha_1,r\alpha_2$, respectively, for every sufficiently small real
$r$ (and by continuation along the leaves for all $r$).
Choose such $r$ with $1,\xi,r$ rationally independent. A rational relation
between the new first coordinates would annihilate both $\beta$ and
$\alpha$, which is impossible. The moved pair has $r_x=2$ and its orbit
closure is $M\times M$. If one original rank is zero, genericity forces
the other to be two, already treated.
\end{proof}

\begin{corollary}\label{cor:relation}
Every surface-dark pair $(s,t)$ satisfies
$a\pi(s)+b\pi(t)\in\Q^2/\Z^2$ for some nonzero integer pair $(a,b)$.
If both torus coordinates are irrational, then $ab\ne0$.
\end{corollary}
\begin{proof}
Surface-lit pairs form a nonempty open invariant set. A dense orbit must
meet that set. If $a=0$ or $b=0$, the corresponding coordinate is rational.
\end{proof}

\section{Rational relation surfaces}
Put $M^\circ=M\setminus L$.
\begin{lemma}[A dark relation component]\label{lem:relation}
Suppose $z=\pi(s)$ is irrational, $a,b$ are relatively prime nonzero
integers, and $az+b\pi(t)=c\in\Q^2/\Z^2$.
Let $F_0$ be the connected component through $(s,t)$ of
\[
 F=\{(s',t')\in M^\circ\times M^\circ:
                   a\pi(s')+b\pi(t')=c\}.
\]
If $(s,t)$ is surface-dark, every distinct pair in $F_0$ is surface-dark.
\end{lemma}
\begin{proof}
Choose rational $z_0,w_0\in\T^2$ with $az_0+bw_0=c$.
Since $\gcd(a,b)=1$, the solution torus downstairs is parametrized
bijectively by
\begin{equation}\label{eq:param}
 z'=z_0+b\zeta,\qquad w'=w_0-a\zeta,
 \qquad \zeta\in\T^2.
\end{equation}
For example, Bezout coefficients give an inverse to the map
$\zeta\mapsto(b\zeta,-a\zeta)$.
Remove the finite rational set $E$ where $z'$ or $w'$ is zero.
The pullback $F$ is a finite covering of $\T^2\setminus E$.
Completing each connected component adds finitely many points and produces
a compact translation surface branched only over rational points.
After subdividing a rational grid, it is tiled by squares of side $1/N$.
The parameter $\zeta$ of $(s,t)$ is irrational because $z$ is irrational.

A power of the diagonal action $h\times h$ preserves the rational relation
and the component $F_0$: the orbit of $c$ on its finite rational torus grid
is finite, and there are finitely many components. In the chart $\zeta$
its action is affine with derivative a power of $H$ and rational
translation part. Taking another power makes this translation part zero
modulo $\Z^2$ and fixes every vertex of a sufficiently fine rational tile
grid upstairs. To justify the first step, the affine action permutes a finite
rational grid containing $0$; a power fixes $0$. To justify the second,
its action on the finite set of preimages of the grid has finite order.
The same construction applies to $v\times v$ and $V$.
When applying Lemma~\ref{lem:one}, take further powers to fix tile edges
pointwise as well. That lemma makes the orbit of $(s,t)$ under the resulting
subgroup dense in the completed component, and hence dense in $F_0$.

If a distinct pair in $F_0$ were surface-lit, its open illumination
neighbourhood would meet this orbit. Affine invariance would make $(s,t)$
surface-lit, a contradiction.
\end{proof}

\section{Proof of the reduction}
\begin{proof}[Proof of Theorem~\ref{thm:slide}]
Write $S=(x,y)$, $T=(u,w)$, $s=(S,(1,1))$, and $t_e=(T,e)$.
All four pairs $(s,t_e)$ are surface-dark by~\eqref{eq:unfold}.
If $T$ were rational, Lemma~\ref{lem:mixed} would contradict darkness.
Thus both $S$ and $T$ are irrational points.

For each $e$ the relation corollary supplies nonzero integers $a_e,b_e$ with
\begin{equation}\label{eq:allrelations}
 a_e(x,y)+b_e(e_1u,e_2w)\in\Q^2.
\end{equation}
Divide by their greatest common divisor to make them relatively prime;
this preserves rationality of the right-hand side.
Suppose both $x$ and $y$ were irrational. From $e=(1,1)$ we obtain
$u=r x$ and $w=r y$ modulo $\Q$, for nonzero rational $r$.
From $e=(-1,1)$ we obtain $u=r'x$ and $w=-r'y$ modulo $\Q$.
Then $(r-r')x$ and $(r+r')y$ are rational, forcing $r=r'$ and $r=-r'$,
a contradiction. Exactly one coordinate of $S$ is irrational.
Exchange axes so that it is $x$.
The relation for $e=(1,1)$ now gives $y,w\in\Q$ and
$u=\rho x+c$, with $\rho\in\Q\setminus\{0\}$ and $c\in\Q$.
Set $v=w$ to match the statement.

For any sheet $e$, the first coordinate of~\eqref{eq:allrelations} implies
$a_e+e_1b_e\rho=0$, since $x$ is irrational. Therefore the path
\[
 \delta\longmapsto(s_\delta,t_{e,\delta})
\]
obtained from the stated room slides keeps its rational relation constant.
For sufficiently small $\delta$ it remains in $M^\circ\times M^\circ$,
with distinct interior room endpoints. It consequently remains in the
same relation component. Lemma~\ref{lem:relation} makes all four lift pairs
surface-dark, so~\eqref{eq:unfold} gives darkness in the room.

For an explicit permissible neighbourhood, take $|\delta|$ smaller than
\[
 \min\left\{\operatorname{dist}(x,\Z),
 \frac{\operatorname{dist}(u,\Z)}{|\rho|},
 \operatorname{dist}(S,\partial P),
 \frac{\operatorname{dist}(T,\partial P)}{|\rho|},
 \frac{|S-T|}{1+|\rho|}\right\}.
\]
Every entry is positive. Choosing $x+\delta$ rational gives rational
coordinates for both endpoints and can be done arbitrarily close to zero.
\end{proof}

\begin{proof}[Proof of Theorem~\ref{thm:reduction}]
If a dark distinct interior pair exists and both endpoints are rational,
it already contradicts the hypothesis. Otherwise name an irrational
endpoint $S$ and apply Theorem~\ref{thm:slide} to obtain a rational dark pair.
\end{proof}

\begin{corollary}[Scaling formulation]\label{cor:scaling}
Fix a class of strictly simple orthogonal lattice rooms closed under
positive integer dilations. If every pair of distinct interior lattice
points illuminates in every room of the class, then every distinct real
interior pair illuminates in every room of the class.
\end{corollary}
\begin{proof}
In a fixed room, any rational pair becomes a lattice pair after a common
integer dilation. Dilation preserves legal billiard paths. Apply
Theorem~\ref{thm:reduction}.
\end{proof}

\section{Scope and verification}
The argument is analytic and concerns all real coordinates. It does not
sample the endpoints, enumerate rooms, or impose a reflection cutoff.
The scaling corollary does not convert an area-bounded census into an
all-real theorem: scaling can exceed every given area bound.
Similarly the proof supplies no universal finite list of directions,
reflection words, or rational denominators.

All convex vertex lifts remain forbidden throughout. Interior lattice lifts
and smooth wall lattice lifts are regular, allowed points. The affine maps
used in the proof fix every lattice lift, so they preserve the complete
forbidden set. Their action is used on the unfolding rather than assumed
to act on room pairs.

The source reduction has independent AI mathematical review. This manuscript
incorporates the review's finite-power corrections and the explicit small-slide
bound. The folder provenance identifies the source and review versions.
No computer calculation is needed to establish the quantified theorem.

\paragraph{Acknowledgments and AI assistance.}
The authors used ChatGPT and Claude to develop and check arguments and
prepare the manuscript. The authors take responsibility for its content.

\begin{thebibliography}{9}
\bibitem{Wolecki}
A.~Wolecki, \emph{Illumination in Rational Billiards},
\texttt{arXiv:1905.09358v3} (2024), Lemma~2.1 and Remark~2.2.
\url{https://arxiv.org/abs/1905.09358}.
\end{thebibliography}
\end{document}
